% sections/05-discovery.tex — organized by autonomy level; evidence-grounded per corpus.
% Provenance: all cited values verified in references/txt/* extracts or local PDFs (CP1-final).
\section{Discovery: From Capability Probes to Production-Oriented Agents}
\label{sec:discovery}

Reading the in-window corpus through the decision-authority test yields a discovery landscape
that separates cleanly into pipeline automation (A1) and task-level agents (A2), with no A0
\emph{system} records: prompted-model capability probes persist only as benchmark subjects
(Sec.~\ref{sec:framework}), not as operational discoverers.

\subsection{A1: Pipeline-Orchestrated Discovery at Ecosystem Scale}
The most extensive automated discovery effort in our corpus is Google's AI-augmented
OSS-Fuzz program. Its LLM executes the first four steps of the developer fuzzing workflow ---
drafting fuzz targets, repairing compilation and runtime faults, running campaigns, and
triaging crashes into driver-vs-project attributions --- while step five (fixing) was
explicitly future work at announcement time and an agent-based architecture was listed as a
roadmap item, anchoring its A1 classification~\cite{ossfuzzlvlup2024}. The reported yield is
material: 26 new vulnerabilities across OSS-Fuzz projects, including CVE-2024-9143 in OpenSSL,
assessed by Google as likely present for two decades and ``wouldn't have been discoverable with
existing fuzz targets written by humans''~\cite{ossfuzzlvlup2024}. Two features matter for our
framework. First, the workflow controller, not the model, fixes the stage sequence --- the
defining A1 property. Second, triage automation remains pre-human-review: the authors state the
goal of reaching confidence ``not requiring human review,'' implying review persisted at
announcement time~\cite{ossfuzzlvlup2024}.

\subsection{A2: Seeded Variant Analysis in Production Code}
Project Zero's Big Sleep provides one of the earliest documented production-oriented A2 discoveries in our corpus: an
exploitable stack buffer underflow in SQLite's \texttt{seriesBestIndex}, triggered through the
\texttt{generate\_series} virtual table under a ROWID constraint, found by variant analysis over
a human-curated pool of recent commits using Gemini~1.5~Pro~\cite{bigsleep2024naptime}. Three
qualifiers discipline any capability claim. The seed pool was human-selected; the run's target
was inspired by an AIxCC finding (Sec.~\ref{sec:coevolution}); and the authors themselves note
the result is experimental and that ``a target-specific fuzzer would be at least as effective''
in general~\cite{bigsleep2024naptime}. The built-in baseline is instructive: AFL failed to
rediscover the bug after 150 CPU-hours even with a corpus pre-seeded with the required keywords,
and the OSS-Fuzz harness lacked the \texttt{generate\_series} extension entirely
--- the comparison points to harness and environment coverage as an important limiting factor; the 150 CPU-hour result does not isolate compute from harness design.
No CVE was assigned: the bug was fixed pre-release~\cite{bigsleep2024naptime}.

\subsection{A2: Entry-Point-Assisted Auditing}
DARKNAVY's Argusee distributes auditing across Manager, Auditor, and Checker agents backed by an
LSP-based code reader, and reports the discovery and analysis of CVE-2025-37891 --- an arbitrary
kernel heap overflow in the Linux USB MIDI2-to-UMP conversion path --- alongside fifteen further
vendor-reported flaws in projects such as GPAC and GIFLIB~\cite{darknavyargusee2025}. Its own
text, however, fixes the autonomy boundary: Argusee ``is not intended to completely replace
manual auditing or to discover vulnerabilities from scratch,'' operating from
\emph{human-supplied entry points}; exploitation of CVE-2025-37891 to root on Arch Linux was
performed by the human team, and Reproducer/Exploit agents are listed as future
work~\cite{darknavyargusee2025}. Argusee also reports perfect scores on CyberSecEval~2
buffer-overflow cases --- a saturation datum we return to in Sec.~\ref{sec:validity}.

\subsection{A2: Black-Box Platform Operation}
XBOW operated as a commercial submission engine across public and private HackerOne programs:
an LLM-assisted targeting layer parsed program policies, scored assets, and deduplicated
infrastructure, while internal ``validators'' re-verified findings before
submission~\cite{xbowtop12025}. The vendor reports ${\sim}1{,}060$ submissions of which 130
were resolved and 303 triaged, with 208 duplicates and 209 informative outcomes --- roughly 39
percent of submissions falling into duplicate/informative categories rather than new/resolved
findings, a metric-definition problem for leaderboard-style counting rather than simple noise
--- including a previously unknown Palo Alto GlobalProtect VPN issue affecting 2{,}000+ hosts,
for which no CVE identifier is given~\cite{xbowtop12025}. The decisive classifier for our
taxonomy is the platform rule: ``our security team reviewed them pre-submission to comply with
HackerOne's policy on automated tools''~\cite{xbowtop12025} --- independent reporting fails,
so XBOW is A2 regardless of marketing language.

\subsection{A2 Under Experiment: Production Networks}
The ARTEMIS study places six agent instances against ten professionals on the same
corporate-style network (${\sim}8{,}000$ hosts, 12 subnets), scoring the first ten hours of
sixteen-hour runs~\cite{artemis2025comparing}. We use it as environment-class E2 evidence that
agents can operate on realistic infrastructure; press-derived superiority percentages are
excluded from this paper because they do not appear in the primary text.

\subsection{Cross-Cutting Comparison}
Table~\ref{tab:disccompare} consolidates the discovery records. Reading down the autonomy
column, the pattern consistent with I1 is cross-sectional rather than longitudinal: every record retains an
identified human locus (seed curation, entry points, mandatory review, target selection), and what
we observe is that human involvement increasingly appears at the boundaries of the autonomous
workflow --- seed selection, target selection, entry-point specification, final reporting --- rather
than on each individual tool action, while evaluation, funding, and discourse increasingly score
systems by their degree of autonomous operation.

\begin{table}[t]
\centering\scriptsize
\caption{Discovery-side comparison across autonomy levels. E-class per protocol v4;
validity flags use the V1--V5 taxonomy of Sec.~\ref{sec:validity}. All production statistics are
vendor-reported unless noted.}
\label{tab:disccompare}
\resizebox{\columnwidth}{!}{%
\begin{tabular}{p{1.55cm}llp{1.5cm}p{1.6cm}p{1.7cm}}
\hline
Record & $A$/E & Human locus & Metric & Baseline & Validity risks\\
\hline
OSS-Fuzz AI \cite{ossfuzzlvlup2024} & A1/E2 & report review & 26 vulns; CVE-2024-9143 &
 prior harnesses & V5 triage counts\\
Big Sleep \cite{bigsleep2024naptime} & A2/E2 & curated seeds & n=1 find + repro &
 AFL 150 CPU-h fail & V1 model preknowledge\\
Argusee \cite{darknavyargusee2025} & A2/E2 & entry points & CVE-2025-37891; 15 claims &
 CSEv2 100\% (sat.) & self-adjudication\\
XBOW \cite{xbowtop12025} & A2/E2 & mandated review & funnel stats; rank & leaderboard &
 unaudited dashboard\\
ARTEMIS \cite{artemis2025comparing} & A2/E2 & comparison arm & first-10h score & professionals &
 single environment\\
\hline
\end{tabular}}
\end{table}
